\documentclass[11pt]{article}

\usepackage[margin=1in]{geometry}
\usepackage{amsmath, amssymb, amsthm, mathtools}
\usepackage{enumitem}
\usepackage{booktabs}
\usepackage{hyperref}
\usepackage{graphicx}

\hypersetup{
  colorlinks=true,
  linkcolor=blue,
  citecolor=blue,
  urlcolor=blue
}

% -----------------------------
% Theorem environments
% -----------------------------
\newtheorem{definition}{Definition}
\newtheorem{proposition}{Proposition}
\newtheorem{remark}{Remark}

% -----------------------------
% Macros
% -----------------------------
\newcommand{\R}{\mathbb{R}}
\newcommand{\Rp}{\R_{+}}
\newcommand{\norm}[1]{\left\lVert #1 \right\rVert}
\newcommand{\inner}[2]{\left\langle #1,#2\right\rangle}

\title{\vspace{-1.0em}
\textbf{Project Proposal: A Candidate, Computable LMS$\to$Bloch-Disk Embedding as a Bridge to Geometric--Quantum Models of Color}\\
\vspace{0.25em}
\small \textit{Course: Mathematical Methods for Image Processing}}

\author{
Guy Singer - 206222093}


\date{\underline{\hspace{4.2cm}}}

\begin{document}
\maketitle

\vspace{-0.5em}
\begin{abstract}
Recent work in mathematical neuroscience develops intrinsic geometric models of perceived color space:
(i) an axiomatic/homogeneous-cone program revisiting Resnikoff's classification and its metric consequences, and
(ii) a Jordan-algebra (``rebit'') reformulation in which chromatic states form a 2D Bloch disk identified with Klein/Hilbert hyperbolic geometry.
While mathematically elegant, these theories are largely intrinsic and do not specify a concrete computational map from physical/device representations (RGB, spectral stimuli, or LMS cone responses) to the abstract state/effect objects.
Moreover, key assumptions about background transformations and metric invariance are explicitly identified as open or empirically problematic.

This proposal is therefore scoped in two stages.
\textbf{Part I (course scope, 2 weeks)} does \emph{not} attempt to identify the ``true'' perceptual embedding.
Instead, it produces a \emph{candidate, computable} embedding pipeline
\[
\text{sRGB} \to \text{linear RGB} \to \text{XYZ} \to (L,M,S) \to (v_1,v_2)\in D \to \rho(v_1,v_2),
\]
motivated by opponent processing and constrained so that outputs lie in the rebit Bloch disk.
We then analyze basic mathematical properties, implement hyperbolic/Hilbert distances on the disk, and demonstrate one image-processing application (e.g.\ geodesic chromatic interpolation).
\textbf{Part II (after the course completion)} outlines a research program aimed at publishable novelty: calibrating and validating the embedding and induced distances on discrimination data; estimating observer-specific convex domains in Hilbert/Finsler geometry (linked to MacAdam ellipses); and extending the framework to context effects (background dependence) via structured state transformations (e.g.\ open-system channels), motivated by known failures of background-invariant Riemannian metrics.
\end{abstract}

\section{Background and Motivation}

\subsection{Axiomatic geometry of perceived color}
Under standard colorimetric conditions, the space $P$ of perceived colors can be modeled as a regular convex cone embedded in a real vector space of dimension at most $3$ (Schr\"odinger/Grassmann/von Helmholtz axioms).
Adding Resnikoff's homogeneity axiom leads to a strong classification: $P$ must be diffeomorphic to either a ``flat'' cone $P_1 \cong \Rp^3$ or a ``hyperbolic'' model $P_2 \cong \Rp \times \mathrm{SL}(2,\R)/\mathrm{SO}(2)$.
Crucially for this proposal, the background-transformation hypothesis requires linearity properties that remain experimentally unverified, and the background-invariance assumption used to single out Riemannian metrics is contradicted by context phenomena (e.g.\ crispening), motivating further work on perceptual metrics and contextual dependence \cite{Provenzi2020}.

\subsection{Jordan algebras and the rebit / Bloch disk}
A complementary approach replaces the external homogeneity postulate with a structural axiom:
the Grassmann--von Helmholtz cone is the positive cone of a formally real Jordan algebra of dimension $3$.
The classification yields either $\R\oplus\R\oplus\R$ (flat case) or $H(2,\R)$ (real symmetric $2\times2$ matrices), corresponding to the hyperbolic case.
In the $H(2,\R)$ case, the chromatic state space is a 2D Bloch disk, isometric to a Klein disk; the induced metric coincides with the Hilbert metric on the disk, whose geodesics are Euclidean chords \cite{Berthier2020}.
This provides a natural intrinsic description of opponency using two Pauli-like generators.

\subsection{The specific gap this project targets}
Both approaches are largely intrinsic:
they define spaces, symmetries, and metrics internally, but do not provide a concrete computational mapping from cone excitations or device RGB to the rebit coordinates $(v_1,v_2)$ (nor claim identifiability of such a map from the axioms alone).
Indeed, the rebit construction is explicitly presented without reference to physical colors or to an observer \cite{Berthier2020}.
This missing bridge motivates the present proposal.

\section{Problem Statement and Explicit Scope}

\subsection{Problem statement}
We aim to build a mathematically controlled and computationally implementable \emph{candidate} map from LMS cone responses (and thus from RGB images via standard conversions) into the rebit Bloch disk coordinates, enabling:
(i) reproducible computations in the geometry proposed by the papers,
and (ii) preliminary algorithmic demonstrations in image processing.

\subsection{Scope and non-goals (important)}
\paragraph{Part I \emph{does not} claim to:}
\begin{itemize}[leftmargin=2.2em]
\item identify the unique or ``true'' human perceptual embedding LMS$\to$Bloch disk;
\item establish empirical validity across observers, viewing conditions, or strong context effects;
\item resolve the unverified linearity of background transformations or crispening-compatible metrics.
\end{itemize}

\paragraph{Part I \emph{does} aim to:}
\begin{itemize}[leftmargin=2.2em]
\item propose a small, explicit family of candidate LMS$\to$disk embeddings motivated by opponent processing;
\item prove/verify basic mathematical properties (disk-boundedness, achromatic mapping, symmetries);
\item implement the geometry (Klein/Hilbert distance) and demonstrate one image-processing use-case.
\end{itemize}

This scoping distinction aligns with the notion that identifying/validating the true mapping is thesis-level work, whereas constructing and analyzing a candidate computable bridge is feasible as a course final project.

\section{Part I: Two-Week Course Project}

\subsection{Inputs, outputs, and assumptions}
\paragraph{Inputs.}
Color images in sRGB (or other known device space), converted to linear RGB and then to XYZ and LMS under a fixed, specified conversion pipeline.

\paragraph{Assumptions.}
We fix a viewing condition and a reference white point (or apply a simple chromatic adaptation step before embedding).
We treat chromaticity and intensity separately to avoid conflating brightness/lightness with chromatic state in the first implementation.

\paragraph{Outputs.}
For each pixel (or sampled color) we compute:
\begin{itemize}[leftmargin=2.2em]
\item LMS cone responses $x=(L,M,S)\in\Rp^3$,
\item Bloch-disk coordinates $v(x)=(v_1(x),v_2(x))$ with $\norm{v(x)}<1$,
\item associated rebit state density matrix $\rho(v_1,v_2)$ as in \cite{Berthier2020}.
\end{itemize}

\subsection{Candidate LMS$\to$Bloch-disk embedding}

\subsubsection{Opponent coordinates}
Let $x=(L,M,S)\in\Rp^3$.
Define a luminance proxy
\begin{equation}
Y(x) := w_L L + w_M M,\qquad w_L,w_M>0,
\end{equation}
and two opponent signals
\begin{equation}
O_1(x) := L - \gamma M,\qquad
O_2(x) := S - \beta(L+M),
\end{equation}
with parameters $\gamma,\beta>0$.
These represent a standard ``two-opponent + luminance'' decomposition consistent with classical pre-cortical opponency discussions and the rebit's two Pauli-like generators \cite{Berthier2020}.

\subsubsection{Luminance-normalized chromatic coordinates}
Normalize by luminance to remove global intensity scaling:
\begin{equation}\label{eq:u}
u(x) := (u_1(x),u_2(x))
= \left(
\frac{O_1(x)}{Y(x)+\varepsilon},\;
\frac{O_2(x)}{Y(x)+\varepsilon}
\right),
\qquad \varepsilon>0.
\end{equation}

\subsubsection{Disk-bounded squashing (guaranteeing $\norm{v}<1$)}
Define a smooth radial squashing map $T:\R^2\to D$ where $D=\{v\in\R^2:\norm{v}<1\}$:
\begin{equation}\label{eq:squash}
v(x) := T(u(x)) :=
\begin{cases}
(0,0), & u(x)=(0,0),\\[0.35em]
\tanh(\kappa\norm{u(x)})\dfrac{u(x)}{\norm{u(x)}}, & \text{otherwise},
\end{cases}
\qquad \kappa>0.
\end{equation}
The parameter $\kappa$ controls how quickly colors saturate toward the boundary.

\subsubsection{Rebit state (Bloch disk)}
Let $\sigma_1,\sigma_2$ be the Pauli-like matrices (real symmetric) used for rebits in \cite{Berthier2020}:
\begin{equation}
\sigma_1=\begin{pmatrix}1&0\\0&-1\end{pmatrix},\qquad
\sigma_2=\begin{pmatrix}0&1\\1&0\end{pmatrix}.
\end{equation}
Define the state density matrix:
\begin{equation}\label{eq:rho}
\rho(x):=\rho(v_1(x),v_2(x))
=\frac12\left(I_2+v_1(x)\sigma_1+v_2(x)\sigma_2\right),
\end{equation}
which is well-defined and positive semidefinite whenever $\norm{v(x)}\le 1$.

\subsubsection{Basic properties to prove/check}
\begin{proposition}[Disk-boundedness]
For all $x\in\Rp^3$, the map \eqref{eq:squash} satisfies $\norm{v(x)}<1$. Hence $\rho(x)$ in \eqref{eq:rho} is a valid rebit state density matrix.
\end{proposition}

\begin{proposition}[Achromatic mapping]
If $O_1(x)=0$ and $O_2(x)=0$ (chromatic equilibrium in both opponencies), then $u(x)=(0,0)$ and therefore $v(x)=(0,0)$ and $\rho(x)=I_2/2$, the maximal-entropy state in \cite{Berthier2020}.
\end{proposition}

\begin{remark}[Status of the construction]
The mapping above is intentionally presented as a \emph{candidate family} with parameters $(w_L,w_M,\gamma,\beta,\kappa)$.
Part I does not claim identifiability or perceptual optimality; it provides an explicit bridge to enable computation in the intrinsic geometry and to set up Part II calibration/validation.
\end{remark}

\subsection{Geometry on the disk and computational implementation}

\subsubsection{Hilbert/Klein distance on the unit disk}
Berthier emphasizes that the Klein disk geometry is well adapted to the rebit state space and that the Klein metric coincides with Hilbert's metric on the disk \cite{Berthier2020}.
We will implement the Hilbert distance $d_H$ between two points $p,q\in D$:
let the Euclidean line through $p,q$ intersect the unit circle at boundary points $r,s$ in the order $r,p,q,s$.
Then
\begin{equation}
d_H(p,q)=\frac12\log[r,p,q,s],\qquad
[r,p,q,s]=\frac{\norm{q-r}}{\norm{p-r}}\cdot\frac{\norm{p-s}}{\norm{q-s}}.
\end{equation}
This gives a closed-form computable distance on the disk (with robust numerical handling for near-boundary cases).

\subsubsection{One image-processing demonstration}
We will implement at least one of the following demonstrations:

\paragraph{Demo A: Hyperbolic/geodesic chromatic interpolation.}
Given two colors $x_0,x_1$ in an image, embed to $v_0,v_1\in D$ and interpolate along a Klein geodesic (Euclidean chord), then map back through an inverse display pipeline (with gamut handling).
Compare against linear interpolation in RGB and in a baseline opponent space.

\paragraph{Demo B: Adaptation-then-embedding.}
Apply a simple von Kries diagonal adaptation in LMS prior to embedding, and show the resulting shift of the embedded chromatic distribution of an image under different assumed illuminants.
This connects computationally to the role of illuminant/background changes emphasized in the homogeneous-cone framework \cite{Provenzi2020}.

\subsection{Deliverables}
\begin{itemize}[leftmargin=2.2em]
\item \textbf{Report (8--12 pages)} with clear assumptions and non-claims, derivation of the embedding family, and computational results.
\item \textbf{Reproducible code} (Python/Jupyter) implementing: RGB$\to$XYZ$\to$LMS, embedding \eqref{eq:u}--\eqref{eq:rho}, Hilbert distance, and one demo.
\item \textbf{Figures} showing: disk embeddings of color ramps and image samples, opponent-axis behavior, and example outputs of the demo.
\end{itemize}

\section{Part II: Post-Course Research Extension (Toward Publication)}

\subsection{Motivation for extension}
Part I produces a candidate computational bridge but does not claim perceptual correctness.
Part II targets scientific novelty by turning the candidate into an empirically constrained and observer-/context-aware model, aligned with open problems emphasized by the papers:
(i) unverified linearity and problematic background-invariant metrics \cite{Provenzi2020},
(ii) the intrinsic nature of the rebit model and the suggestion to relate Hilbert/Finsler geometry to MacAdam ellipses \cite{Berthier2020},
and (iii) context effects potentially requiring open-system (channel) descriptions \cite{Berthier2020}.

\subsection{Research questions}
\begin{enumerate}[label=\textbf{RQ\arabic*.}, leftmargin=2.2em]
\item \textbf{Calibration:} Can parameters of the candidate embedding family be fit so that disk-geometry distances predict discrimination data better than baselines?
\item \textbf{Observer geometry:} Can an observer be modeled by a convex subset $\Omega\subset D$ such that the induced Hilbert/Finsler metric locally matches measured MacAdam ellipses, as suggested in \cite{Berthier2020}?
\item \textbf{Context dependence:} Can we model background/context changes by structured transformations on states that are \emph{not} globally isometric (consistent with crispening), while retaining mathematical control?
\end{enumerate}

\subsection{Work packages}

\paragraph{WP-A: Data-driven calibration of the embedding.}
Introduce parameters $\theta=(w_L,w_M,\gamma,\beta,\kappa,\dots)$ and fit them on discrimination datasets (e.g.\ MacAdam ellipses or modern color-difference datasets under controlled conditions).
Compare predictive performance of:
\begin{itemize}[leftmargin=2.2em]
\item Euclidean baselines (Lab, opponent spaces),
\item Klein/Hilbert distances on the embedded disk,
\item variants incorporating nonlinear compression (e.g.\ log-like response) prior to \eqref{eq:u}.
\end{itemize}

\paragraph{WP-B: Estimating an observer-specific convex domain $\Omega$ (Hilbert/Finsler).}
Berthier proposes identifying for each observer a convex subset $\Omega$ by comparing Finsler metric balls to MacAdam ellipses \cite{Berthier2020}.
We formalize this as an inverse problem:
parameterize $\partial\Omega$ (support function, spline, or convex polygon) and optimize $\Omega$ so that the local unit balls of the Hilbert/Finsler norm match measured ellipses (up to scale).
This yields a principled, observer-specific geometry.

\paragraph{WP-C: Context effects via structured state maps (open-system/channel model).}
Provenzi shows background-invariant Riemannian metrics conflict with crispening and highlights the need for further research on perceptual metrics in context \cite{Provenzi2020}.
Berthier notes that context effects may be modeled by open quantum systems via linear, trace-preserving, completely positive maps \cite{Berthier2020}.
We propose a family of context-indexed maps $\mathcal{E}_b$ acting on states:
\[
\rho \mapsto \mathcal{E}_b(\rho),
\]
and define a context-dependent distance
\[
d_b(x,y):=d\!\left(\mathcal{E}_b(\rho(x)),\mathcal{E}_b(\rho(y))\right),
\]
with the goal of reproducing crispening-like modulations while preserving tractable geometry.

\paragraph{WP-D (optional): Empirical test of additivity/linearity for background transformations.}
If feasible, implement (or simulate) the psycho-physical protocol proposed in \cite{Provenzi2020} to test additivity of background transformations.
This provides direct evidence about a core open hypothesis in the homogeneous-cone approach.

\subsection{Expected novelty and publication plan}
The publishable contribution is not merely ``a mapping,'' but:
\begin{enumerate}[leftmargin=2.2em]
\item a calibrated and validated embedding with quantitative predictive evaluation;
\item an algorithm to infer observer-specific convex domains $\Omega$ with Hilbert/Finsler geometry linked to discrimination ellipses;
\item a mathematically controlled context-dependent extension using structured state maps, directly motivated by the failures of background-invariant metrics.
\end{enumerate}
Potential venues include Journal of Mathematical Neuroscience, SIAM Journal on Imaging Sciences, and interdisciplinary vision-science venues depending on the empirical component.

\section{Risks and Mitigations}
\begin{itemize}[leftmargin=2.2em]
\item \textbf{Underdetermined mapping:} Part I embraces non-uniqueness; Part II addresses it via calibration (WP-A) and cross-validation.
\item \textbf{Dataset mismatch / viewing conditions:} Explicitly fix assumptions; use datasets with documented viewing geometry; report sensitivity analyses.
\item \textbf{Context complexity:} Start with controlled backgrounds; treat open-system/channel models as incremental extensions (WP-C).
\item \textbf{Human-subject constraints:} Keep Part I non-human-subject; only pursue WP-D with appropriate approvals if needed.
\end{itemize}

\section*{Approval Request}
I request approval to pursue:
\begin{enumerate}[leftmargin=2.2em]
\item \textbf{Part I (2-week final project):} construct and implement a candidate, computable LMS$\to$Bloch-disk embedding family; implement Hilbert/Klein distance on the disk; and demonstrate one image-processing application.
\item \textbf{Part II (post-course):} develop a publishable extension by calibrating/validating the embedding, fitting observer-specific convex domains via Hilbert/Finsler geometry, and extending to context effects via structured state maps.
\end{enumerate}

\vspace{0.5em}
\bibliographystyle{plain}
\begin{thebibliography}{9}

\bibitem{Provenzi2020}
E.~Provenzi.
\newblock \emph{Geometry of color perception. Part 1: structures and metrics of a homogeneous color space}.
\newblock Journal of Mathematical Neuroscience, 10:7, 2020.
\newblock DOI: 10.1186/s13408-020-00084-x.

\bibitem{Berthier2020}
M.~Berthier.
\newblock \emph{Geometry of color perception. Part 2: perceived colors from real quantum states and Hering's rebit}.
\newblock Journal of Mathematical Neuroscience, 10:14, 2020.
\newblock DOI: 10.1186/s13408-020-00092-x.

\end{thebibliography}

\end{document}

